\documentclass[11pt]{article}

\usepackage[margin=1in]{geometry}
\usepackage{amsmath,amssymb,amsthm,mathtools}
\usepackage{graphicx}
\usepackage{microtype}
\usepackage{enumitem}
\usepackage[hidelinks]{hyperref}

\newtheorem{theorem}{Theorem}
\newtheorem{proposition}{Proposition}
\newtheorem{remark}{Remark}
\newcommand{\sot}{\mathrm{SOT}}
\newcommand{\1}{\mathbf 1}

\title{Uniqueness Without Monotonicity:\\
Scalar Dynamics for Spectral Transforms}
\author{Candidate full answer to Appendix A, Problem 1 of arXiv:2508.04890}
\date{17 August 2026}

\begin{document}
\maketitle

\begin{center}
\fbox{\parbox{0.92\linewidth}{\textbf{Status.} Candidate full answer,
likely valid, subject to human review.  The source's phrase ``order of
spectral projections'' is not defined; the result treats both the standard
spectral/Loewner order and the source's formally stated projection-stability
axiom.}}
\end{center}

\section{Source problem}

The target is Appendix A, Problem 1 (PDF p.~11) in Alpay--Alakkad--Alpay,
\emph{Transfinite Operator Fixed Points on Hilbert Spaces: An Alpay Algebra
Approach}, arXiv:2508.04890v1 \cite{source}.  It asks whether the monotonicity
assumption is essential, requests a non-order-preserving transformation with
a unique iterated limit, and alternatively asks for weaker hypotheses that
guarantee uniqueness.

\begin{figure}[ht]
  \centering
  \includegraphics[width=\textwidth]{figures/open_problem_crop.png}
  \caption{Appendix A, Problem 1 in the source paper (PDF p.~11).}
\end{figure}

The source's formal definition uses a stability condition that, in
particular, requires \(\Phi(P)=P\) for every orthogonal projection \(P\).
Its appendix instead refers to preservation of an undefined ``order of
spectral projections.''  We resolve both natural readings below.

\section{Definitions and result}

Fix a nonzero complex Hilbert space \(H\), a compact set
\(K\subset\mathbb R\), and
\[
 \mathcal S_K(H)=\{A\in B(H):A=A^*,\ \sigma(A)\subseteq K\}.
\]
For \(f\in C(K,K)\), let \(\Phi_f(A)=f(A)\), using the continuous functional
calculus.  Write \(f^{\circ n}\) for scalar composition.  For reference, the
spectral order may be defined by
\[
 A\preccurlyeq_s B
 \quad\Longleftrightarrow\quad
 E_A((t,\infty))\le E_B((t,\infty))\quad(t\in\mathbb R).
\]
It agrees with the ordinary order on scalar multiples of the identity.

\subsection*{Proof intuition}

A spectral transform does not need an order structure to converge.  Every
operator orbit is obtained by running the same one-dimensional dynamical
system \(t\mapsto f(t)\) at each spectral value.  If the scalar orbits
converge, the spectral theorem integrates those pointwise limits; bounded
convergence turns them into strong operator convergence.  Thus uniqueness is
controlled by scalar attraction, not by monotonicity.  Squaring on
\([-1,1]\) is already nonmonotone but converges pointwise to the indicator of
\(\{-1,1\}\).  Damping the square by \(1/2\) makes zero a uniform global
attractor and even destroys the source's projection stability.

\begin{theorem}[Scalar-orbit criterion]\label{thm:criterion}
Let \(K\subset\mathbb R\) be compact, \(f\in C(K,K)\), and suppose that
\[
 f^{\circ n}(t)\longrightarrow g(t)\qquad(t\in K).
\]
The pointwise limit \(g\) is Borel, so \(g(A)\) is defined by the bounded
Borel functional calculus.
Then, for every \(A\in\mathcal S_K(H)\),
\[
 \Phi_f^n(A)=f^{\circ n}(A)\xrightarrow[n\to\infty]{\sot}g(A),
 \qquad \Phi_f(g(A))=g(A).
\]
Consequently the orbit limit is unique and is reached at the first limit
stage \(\omega\); no order-monotonicity assumption is needed.  If the scalar
convergence is uniform on \(K\), the operator convergence is in norm.

If all scalar orbits have one common limit \(\xi\in K\), then
\(\Phi_f^n(A)\to\xi I\) strongly for every \(A\), and \(\xi I\) is the unique
fixed point of \(\Phi_f\) in \(\mathcal S_K(H)\).  Conversely, convergence to
the same scalar operator \(\xi I\) for every \(A\) forces
\(f^{\circ n}(t)\to\xi\) for every \(t\in K\).
\end{theorem}

\begin{proof}
Functoriality of the continuous functional calculus gives
\(\Phi_f^n(A)=f^{\circ n}(A)\).  Let \(E_A\) be the spectral measure of \(A\)
and \(\mu_x(\Delta)=\langle E_A(\Delta)x,x\rangle\).  Since every iterate and
the pointwise limit take values in the bounded set \(K\), dominated
convergence gives, for every \(x\in H\),
\[
 \bigl\|(f^{\circ n}(A)-g(A))x\bigr\|^2
 =\int_K |f^{\circ n}(t)-g(t)|^2\,d\mu_x(t)
 \longrightarrow0.
\]
This is strong convergence.  If the convergence on \(K\) is uniform, then
\[
 \|f^{\circ n}(A)-g(A)\|
 \le \sup_{t\in K}|f^{\circ n}(t)-g(t)|\longrightarrow0.
\]

For each \(t\), continuity of \(f\) yields
\[
 f(g(t))=\lim_n f(f^{\circ n}(t))
 =\lim_n f^{\circ(n+1)}(t)=g(t).
\]
Hence \(f\circ g=g\), and functional calculus gives
\(\Phi_f(g(A))=g(A)\).

Now suppose \(g\equiv\xi\).  The preceding argument gives the common limit
\(\xi I\).  If \(B=f(B)\), then the continuous function \(h(t)=t-f(t)\)
satisfies \(h(B)=0\), so \(h\) vanishes on \(\sigma(B)\).  Every
\(t\in\sigma(B)\) is therefore a scalar fixed point.  Its scalar orbit is
constant and also must converge to \(\xi\), whence \(t=\xi\).  Thus
\(\sigma(B)\subseteq\{\xi\}\) and \(B=\xi I\).  Conversely, apply the claimed
common operator convergence to the scalar operator \(A=tI\).
\end{proof}

\section{Two explicit answers to Problem 1}

\begin{proposition}[The source's canonical squaring map]\label{prop:square}
On \(\mathcal S_{[-1,1]}(H)\), define \(\Phi(A)=A^2\).  This transformation
is neither Loewner-order nor spectral-order preserving, yet every orbit has
the unique strong limit
\[
 \Phi^n(A)=A^{2^n}\xrightarrow[n\to\infty]{\sot}
 E_A(\{-1,1\}).
\]
The limit is an orthogonal projection and hence is fixed by \(\Phi\).
\end{proposition}

\begin{proof}
For \(f(t)=t^2\) on \([-1,1]\),
\[
 f^{\circ n}(t)=t^{2^n}\longrightarrow
 g(t)=\1_{\{-1,1\}}(t).
\]
Theorem~\ref{thm:criterion} gives the stated strong limit.  It is fixed
because its square equals itself.  Nonmonotonicity already occurs on scalar
operators: \(-I\le0\) and \(-I\preccurlyeq_s0\), but
\(\Phi(-I)=I\) is neither \(\le0=\Phi(0)\) nor
\(\preccurlyeq_s0\).
\end{proof}

This also corrects the source's motivating claim: for self-adjoint
contractions the limit is generally \(E_A(\{-1,1\})\), not only the
projection onto \(\ker(I-A)\).  The one-dimensional case \(A=-I\) makes the
difference immediate.

\begin{proposition}[Failure of projection stability with a unique global
fixed point]\label{prop:damped}
On \(\mathcal S_{[-1,1]}(H)\), define
\(\Psi(A)=\tfrac12A^2\).  Then \(\Psi\) is neither Loewner-order nor
spectral-order preserving and violates the source's formal stability axiom:
\(\Psi(P)=P/2\ne P\) for every nonzero orthogonal projection \(P\).
Nevertheless
\[
 \Psi^n(A)=2^{1-2^n}A^{2^n},
 \qquad
 \|\Psi^n(A)\|\le 2^{1-2^n}\longrightarrow0.
\]
Thus \(0\) is the norm limit of every orbit and is the unique fixed point in
the self-adjoint unit ball.
\end{proposition}

\begin{proof}
The scalar identity follows by induction from \(h(t)=t^2/2\), and the norm
estimate uses \(\|A\|\le1\).  If \(B=\Psi(B)\) and \(r=\|B\|\), then
\(r\le r^2/2\).  For \(0<r\le1\) this is impossible, so \(r=0\) and \(B=0\).
The projection calculation is explicit.  Finally, \(-I\le0\) (also in
spectral order), while \(\Psi(-I)=I/2\) lies above \(\Psi(0)=0\), proving
both order failures.
\end{proof}

Together Propositions~\ref{prop:square} and~\ref{prop:damped} answer both
parts of Problem 1: they provide explicit non-order-preserving transforms
with unique iterated limits, while Theorem~\ref{thm:criterion} supplies a
weaker, purely dynamical condition guaranteeing uniqueness.

\section{Verification and limitations}

The proof is dimension-free and includes continuous spectrum; it does not
assume diagonalizability.  The included script
\texttt{code/check\_examples.py} checks the finite-spectrum limits, the exact
damped-iterate formula through seven stages, and the scalar order failure.
It is a sanity check, not evidence replacing the proof.

There are three scope qualifications.
\begin{enumerate}[leftmargin=*]
\item The source never defines ``order of spectral projections.''  We use the
standard spectral order and separately violate its formally stated
projection-stability condition.
\item ``Unique \(\Phi^\infty(A)\)'' is read as a uniquely determined orbit
limit for each initial \(A\).  Proposition~\ref{prop:damped} also gives the
stronger property of a unique global fixed point.
\item The result is in the fixed-H, bounded self-adjoint spectral-transform
setting.  This is a special case of the source's framework and suffices for
the requested construction; it does not repair the paper's broader
transfinite claims.
\end{enumerate}

\paragraph{Novelty search.}
On 17 August 2026, bounded searches covered the run's registry, solution,
attempt, and proof-gap indexes, plus arXiv searches by exact title, arXiv id,
Problem 1 label, the quoted order phrase, and close uniqueness/spectral
keywords.  Only the source paper was located; no later explicit answer or
correction was found.  Because the proof is an elementary application of the
spectral theorem, novelty confidence is moderate rather than high.

\paragraph{Human-review recommendation.}
Check the intended meanings of the source's undefined order phrase and
``unique limit,'' then verify the common-attractor fixed-point argument in
Theorem~\ref{thm:criterion}.  The mathematical computations themselves are
short and should be reproducible directly from the spectral theorem.

\begin{thebibliography}{9}
\bibitem{source}
F. Alpay, H. Alakkad, and T. Alpay,
\emph{Transfinite Operator Fixed Points on Hilbert Spaces: An Alpay Algebra
Approach}, arXiv:2508.04890v1 (2025), Appendix A, Problem 1.
\end{thebibliography}

\end{document}
